\setcounter{chapter}{2}
\chapter{Lebesgue外測度}\label{chap:lebesgue-outer}

\paragraph{本章の目標}
Lebesgue外測度，可算劣加法性

第2章で定義したJordan外測度
\[
m_J^*(A) = \inf\{m_J(E) : A \subset E,\ E \text{ は基本集合}\}
\]
は有限個の半開区間の合併（基本集合）で $A$ を覆い，
体積の下限を取るものであった．
この有限被覆の制約により，
$\QQ \cap [0,1]$ のように稠密に散らばった集合に対しては
$m_J^*(\QQ \cap [0,1]) = 1$ となり，
「薄い」集合であることを検出できなかった．

Lebesgue外測度は，被覆に使える区間の個数を
可算無限個に拡張する：
\[
m^*(A)
= \inf\!\left\{
\sum_{j=1}^{\infty} |Q_j|
: A \subset \bigcup_{j=1}^{\infty} Q_j,\
Q_j \text{ は半開区間}
\right\}.
\]
被覆の自由度が増えた分，より精密な評価が可能になり，
$m^*(\QQ \cap [0,1]) = 0$ が示せるようになる．
代償として有限加法性は失われるが，
代わりに\textbf{可算劣加法性}
（可算個の和で上から抑えられる）が成り立つ．

% ==================================================================
\section{Lebesgue外測度の定義}\label{sec:lebesgue-outer-def}
% ==================================================================

\begin{definition}[可算被覆]\label{def:countable-cover}
集合 $A \subset \RR^d$ に対し，
半開区間の列 $\{Q_j\}_{j=1}^{\infty}$ が
$A$ の\textbf{可算被覆}であるとは
\[
A \subset \bigcup_{j=1}^{\infty} Q_j
\]
が成り立つことをいう．
\end{definition}

\begin{remark}
有限個の区間 $Q_1,\ldots,Q_n$ による被覆は，
$Q_j = \emptyset$（$j > n$）とすれば
可算被覆の特別な場合とみなせる．
また $\RR^d = \bigcup_{n=1}^{\infty} [-n,n)^d$ だから，
任意の集合は可算被覆をもつ．
\end{remark}

\begin{definition}[Lebesgue外測度]\label{def:lebesgue-outer-measure}
集合 $A \subset \RR^d$ の\textbf{Lebesgue外測度}を
\[
m^*(A)
:= \inf\!\left\{
\sum_{j=1}^{\infty} |Q_j|
: \{Q_j\}_{j=1}^{\infty} \text{ は } A \text{ の可算被覆}
\right\}
\]
で定める．値は $[0,\infty]$ を取りうる．
\end{definition}

\begin{remark}\label{rem:jordan-vs-lebesgue-outer}
Jordan外測度との対比を整理する．
\begin{center}
\begin{tabular}{l|l|l}
 & Jordan外測度 $m_J^*$ & Lebesgue外測度 $m^*$ \\
\hline
被覆 & 有限個の区間の合併（基本集合） & 可算個の区間 \\
加法性 & 有限加法的 & 可算劣加法的 \\
定義域 & 有界集合 & すべての集合
\end{tabular}
\end{center}
可算被覆は有限被覆を含むから，
下限の候補が増え，
$m^*(A) \le m_J^*(A)$ が常に成り立つ．
\end{remark}

% ==================================================================
\section{基本性質}\label{sec:lebesgue-outer-prop}
% ==================================================================

\begin{theorem}[Lebesgue外測度の基本性質]\label{thm:lebesgue-outer-properties}
\begin{enumerate}
\item （非負性）任意の $A \subset \RR^d$ に対し
$m^*(A) \ge 0$．また $m^*(\emptyset) = 0$．
\item （単調性）$A \subset B \Rightarrow m^*(A) \le m^*(B)$．
\item （可算劣加法性）
\[
m^*\!\left(\bigcup_{n=1}^{\infty} A_n\right)
\le \sum_{n=1}^{\infty} m^*(A_n).
\]
\item （平行移動不変性）
$c \in \RR^d$ に対し $m^*(A+c) = m^*(A)$．
\end{enumerate}
\end{theorem}
\begin{proof}
\textbf{(1)}\quad
体積は非負だから $m^*(A) \ge 0$．
$\emptyset$ は $Q_j = \emptyset$ で覆えるので $m^*(\emptyset) = 0$．

\textbf{(2)}\quad
$B$ の任意の可算被覆は $A$ の可算被覆でもある．
下限の候補が増えるから $m^*(A) \le m^*(B)$．

\textbf{(3)}\quad
いずれかの $m^*(A_n) = \infty$ なら右辺 $= \infty$ で自明．
すべて有限とする．
$\eps > 0$ を固定し，各 $n$ に対して
$A_n$ の可算被覆 $\{Q_{n,j}\}_{j=1}^{\infty}$ を
\[
\sum_{j=1}^{\infty} |Q_{n,j}| < m^*(A_n) + \frac{\eps}{2^n}
\]
となるように取る（下限の定義より可能）．
$\{Q_{n,j}\}_{n,j=1}^{\infty}$ は
$\bigcup_n A_n$ の可算被覆だから
\[
m^*\!\left(\bigcup_{n=1}^{\infty} A_n\right)
\le
\sum_{n=1}^{\infty}\sum_{j=1}^{\infty} |Q_{n,j}|
<
\sum_{n=1}^{\infty}\!\left(m^*(A_n) + \frac{\eps}{2^n}\right)
=
\sum_{n=1}^{\infty} m^*(A_n) + \eps.
\]
$\eps > 0$ は任意だったので結論を得る．

\textbf{(4)}\quad
$\{Q_j\}$ が $A$ の被覆ならば $\{Q_j + c\}$ は $A+c$ の被覆であり，
$|Q_j + c| = |Q_j|$．逆も同様であるから $m^*(A+c) = m^*(A)$．
\end{proof}

\begin{remark}\label{rem:subadditivity-technique}
(3) の証明で用いた
「各 $A_n$ の被覆に $\eps/2^n$ の余裕を持たせて足し上げる」
という手法は，Lebesgue測度論の随所に現れる基本的な技法である．

可算劣加法性は $A_j = \emptyset$（$j > N$）とすれば
有限劣加法性を含む．
一般には等号は成り立たない（加法性ではなく劣加法性）．
加法性を回復するには，
適切な「可測集合」のクラスに制限する必要がある（第5章以降）．
\end{remark}

ここまでの性質は外測度の定義から比較的容易に従う．
次の定理は非自明であり，
Lebesgue外測度が区間の体積と整合することを保証する．

\begin{theorem}[区間の外測度]\label{thm:outer-measure-rectangle}
半開区間 $Q = \prod_{k=1}^d [a_k, b_k)$ に対して
\[
m^*(Q) = |Q| = \prod_{k=1}^d (b_k - a_k).
\]
\end{theorem}
\begin{proof}
\textbf{上界}\quad
$Q$ 自身（と空集合の列）が $Q$ の可算被覆だから
$m^*(Q) \le |Q|$．

\textbf{下界}\quad
$Q$ の任意の可算被覆 $\{Q_j\}_{j=1}^{\infty}$ に対して
$\sum_{j=1}^{\infty} |Q_j| \ge |Q|$ を示す．

$\eps > 0$ を取る．$Q$ をわずかに縮めた閉区間
\[
\overline{Q}_\delta
= \prod_{k=1}^d [a_k + \delta,\ b_k - \delta]
\qquad (\delta > 0 \text{ は十分小さく取る})
\]
を考え，各 $Q_j$ をわずかに膨らませた開区間 $Q_j^\circ$ を
\[
|Q_j^\circ| < |Q_j| + \frac{\eps}{2^j}
\]
となるように取る（各辺を少しずつ広げればよい）．
すると
$\overline{Q}_\delta \subset Q
\subset \bigcup_j Q_j
\subset \bigcup_j Q_j^\circ$
であり，$\overline{Q}_\delta$ はコンパクト（有界閉集合）だから，
\textbf{Heine--Borelの被覆定理}により
有限部分被覆が取れる：
\[
\overline{Q}_\delta
\subset Q_{j_1}^\circ \cup \cdots \cup Q_{j_N}^\circ.
\]
第2章の基本集合に対する有限劣加法性（Jordan外測度の単調性）より
\[
|\overline{Q}_\delta|
\le \sum_{i=1}^N |Q_{j_i}^\circ|
\le \sum_{j=1}^{\infty} |Q_j^\circ|
< \sum_{j=1}^{\infty} |Q_j| + \eps.
\]
$\delta \to 0$ で $|\overline{Q}_\delta| \to |Q|$ だから
$|Q| \le \sum_{j=1}^{\infty} |Q_j| + \eps$．
$\eps > 0$ は任意だったので
$|Q| \le \sum_{j=1}^{\infty} |Q_j|$．

以上の被覆が任意だったので $m^*(Q) \ge |Q|$．
\end{proof}

\begin{remark}
下界の証明で\textbf{Heine--Borelの定理}（コンパクト性）が
本質的に使われている．
半開区間を閉区間で内側から近似し（コンパクト性を確保），
被覆の区間を開区間で外側から近似して（開被覆を構成），
有限部分被覆に帰着させる．
この「$\eps$ だけ余裕を持たせる」技法は
可算と有限の橋渡しとして測度論の基本的な道具である．
\end{remark}

\begin{lemma}[基本集合に対する加法性]\label{lem:outer-additive-elementary}
基本集合 $E$ と任意の集合 $S \subset \RR^d$ が互いに素
（$E \cap S = \emptyset$）ならば
\[
m^*(E \sqcup S) = m^*(E) + m^*(S) = m_J(E) + m^*(S).
\]
\end{lemma}
\begin{proof}
$m^*(E \sqcup S) \le m^*(E) + m^*(S)$
は可算劣加法性から従う．

逆に $E \sqcup S$ の任意の可算被覆 $\{R_k\}$ を取る．
$E = \bigsqcup_{l=1}^n P_l$（$P_l$ は半開区間）と書く．
各 $R_k$ と $P_l$ の交差 $R_k \cap P_l$ は半開区間であり，
$\{R_k \cap P_l\}_k$ は $P_l$ の可算被覆を与える．
定理\ref{thm:outer-measure-rectangle}より
$|P_l| \le \sum_k |R_k \cap P_l|$．

$P_1,\ldots,P_n$ は互いに素で
$\bigsqcup_l (R_k \cap P_l) \subset R_k$ だから
\[
\sum_l |R_k \cap P_l| \le |R_k|.
\]
したがって
\[
m_J(E) = \sum_l |P_l|
\le \sum_l \sum_k |R_k \cap P_l|
= \sum_k \sum_l |R_k \cap P_l|
\le \sum_k |R_k|.
\]
同様に $\{R_k \cap (E \sqcup S) \setminus E\}_k = \{R_k \setminus E\}_k$
は $S$ を覆うから $m^*(S) \le \sum_k |R_k \setminus E|$．
$|R_k| \ge |R_k \cap E| + |R_k \setminus E|$
（ただし $R_k \cap E$ は基本集合で，
$|R_k \cap E| = \sum_l |R_k \cap P_l|$）であるから，
結局
\[
\sum_k |R_k| \ge m_J(E) + m^*(S).
\]
被覆 $\{R_k\}$ は任意だったので
$m^*(E \sqcup S) \ge m_J(E) + m^*(S)$．
\end{proof}

\begin{corollary}\label{cor:outer-measure-elementary}
基本集合 $E = \bigsqcup_{j=1}^n Q_j$ に対して
\[
m^*(E) = m_J(E) = \sum_{j=1}^n |Q_j|.
\]
\end{corollary}
\begin{proof}
$E$ を被覆 $\{Q_1,\ldots,Q_n,\emptyset,\emptyset,\ldots\}$ で覆えば
$m^*(E) \le m_J(E)$．
逆は補題\ref{lem:outer-additive-elementary}を
$E$ 自身と $S = \emptyset$ に適用して
$m^*(E) = m_J(E) + 0 = m_J(E)$．
（あるいは定理\ref{thm:outer-measure-rectangle}の証明と同様の
Heine--Borelの議論を $\overline{E}$ に適用してもよい．）
\end{proof}

% ==================================================================
\section{具体例}\label{sec:lebesgue-outer-examples}
% ==================================================================

\begin{example}[有限集合]\label{ex:finite-null}
有限集合 $A = \{a_1,\ldots,a_n\} \subset \RR^d$ に対して $m^*(A) = 0$．
\end{example}
\begin{proof}
$\eps > 0$ を取る．
各 $a_j$ を一辺の長さ $(\eps/n)^{1/d}$ の半開区間 $Q_j$ で覆えば
$|Q_j| = \eps/n$ であり，$\sum_{j=1}^n |Q_j| = \eps$．
$\eps$ は任意だったので $m^*(A) = 0$．
\end{proof}

\begin{example}[可算集合]\label{ex:countable-null}
可算集合 $A = \{a_1, a_2, \ldots\} \subset \RR^d$ に対して $m^*(A) = 0$．
\end{example}
\begin{proof}
$\eps > 0$ を固定する．
各 $a_n$ を一辺の長さ $(\eps/2^n)^{1/d}$ の半開区間 $Q_n$ で覆えば
$|Q_n| = \eps/2^n$ であり
\[
\sum_{n=1}^{\infty} |Q_n| = \sum_{n=1}^{\infty} \frac{\eps}{2^n} = \eps.
\]
$\{Q_n\}$ は $A$ の可算被覆だから $m^*(A) \le \eps$．
$\eps$ は任意なので $m^*(A) = 0$．
\end{proof}

\begin{example}[{$\QQ \cap [0,1]$ の外測度}]\label{ex:rationals-outer}
$\QQ$ は可算集合であるから，
$\QQ \cap [0,1] \subset \QQ$ の可算劣加法性と単調性より
\[
m^*(\QQ \cap [0,1]) = 0.
\]
一方 Jordan外測度は $m_J^*(\QQ \cap [0,1]) = 1$ であった（第2章）．
Lebesgue外測度は
「有理数の集合は $[0,1]$ の中で薄い」ことを正しく検出している．
\end{example}

\begin{example}[Cantor集合]\label{ex:cantor-outer}
Cantor集合 $C$ を次のように構成する．
\begin{itemize}
\item $C_0 = [0,1]$．
\item $C_1 = [0,\frac{1}{3}] \cup [\frac{2}{3},1]$：
$C_0$ の中央 $\frac{1}{3}$ を除去．
\item $C_n$：$C_{n-1}$ の各閉区間から中央 $\frac{1}{3}$ を除去．
$2^n$ 個の閉区間で，各長さ $3^{-n}$．
\item $C = \bigcap_{n=0}^{\infty} C_n$（Cantor集合）．
\end{itemize}
$C \subset C_n$ であり，
$C_n$ は $2^n$ 個の長さ $3^{-n}$ の区間で覆えるから
\[
m^*(C) \le m^*(C_n) \le 2^n \cdot 3^{-n} = \left(\frac{2}{3}\right)^n.
\]
$n \to \infty$ で $(2/3)^n \to 0$ だから $m^*(C) = 0$．

Cantor集合は\textbf{非可算}
（3進展開により $[0,1]$ と同じ濃度をもつ）にもかかわらず
外測度が $0$ である．
可算集合でなくとも外測度 $0$ になりうることの顕著な例である．
\end{example}

\begin{example}[区間の外測度]\label{ex:interval-outer}
任意の有界区間 $I \subset \RR$（開・閉・半開を問わず）の外測度は
$m^*(I) = |I|$（長さ）に一致する．
\end{example}
\begin{proof}
半開区間 $[a,b)$ については
定理\ref{thm:outer-measure-rectangle}より
$m^*([a,b)) = b - a$．

閉区間：$[a,b) \subset [a,b] \subset [a,b+\eps)$ より
$b - a \le m^*([a,b]) \le b - a + \eps$．
$\eps \to 0$ で $m^*([a,b]) = b - a$．

開区間も同様に
$(a,b) \subset [a,b)$ と $[a-\eps,b) \subset$ を組み合わせて
$m^*((a,b)) = b - a$．
\end{proof}

% ==================================================================
\section{Jordan外測度との比較}\label{sec:jordan-lebesgue-outer}
% ==================================================================

第2章のJordan外測度と本章のLebesgue外測度の関係を
まとめておく．

\begin{theorem}\label{thm:jordan-lebesgue-outer}
任意の $A \subset \RR^d$ に対して $m^*(A) \le m_J^*(A)$．
\end{theorem}
\begin{proof}
$A$ を含む基本集合 $E$ は $A$ の（有限）可算被覆を与える．
系\ref{cor:outer-measure-elementary}より
$m^*(A) \le m^*(E) = m_J(E)$．
$E$ について下限を取れば $m^*(A) \le m_J^*(A)$．
\end{proof}

\begin{example}\label{ex:strict-inequality}
等号が成り立たない例は既にある：
$m^*(\QQ \cap [0,1]) = 0 < 1 = m_J^*(\QQ \cap [0,1])$．
一般に可算集合は常に $m^* = 0$ であるが，
稠密であれば $m_J^* > 0$ を取りうる．
\end{example}

\begin{remark}
一方，Jordan外測度が有限な集合については
$m^*(A) \le m_J^*(A) < \infty$ であり，
特に有界な基本集合については
$m^*(E) = m_J^*(E) = m_J(E)$ が成り立つ（系\ref{cor:outer-measure-elementary}）．
すなわち，Lebesgue外測度はJordan外測度の「改良版」であり，
基本集合の上では一致しつつ，
より広いクラスの集合に対して精密な値を与える．
\end{remark}

\begin{remark}[外測度の限界]\label{rem:outer-limitation}
Lebesgue外測度 $m^*$ は $\calP(\RR^d)$ 全体で定義されるが，
一般には\textbf{加法性}が成り立たない：
互いに素な $A, B$ に対して
$m^*(A \sqcup B) \le m^*(A) + m^*(B)$
は成り立つ（劣加法性）が，逆向きの不等式
$m^*(A \sqcup B) \ge m^*(A) + m^*(B)$
は一般には成り立たない．
（非可測集合の存在は第5章で示す．）

補題\ref{lem:outer-additive-elementary}が示すように，
一方が基本集合であれば加法性が成り立つ．
次章以降では，加法性が回復する集合のクラス
（Lebesgue可測集合）を特定する．
\end{remark}
